\chapter{The Kalman Filter}
\label{ch:kf}

\begin{objectives}
\begin{itemize}
  \item Derive the Kalman filter from Bayes' rule and Gaussian conditioning, in a few lines.
  \item Understand the gain, the innovation and its covariance, and the several equivalent forms of
        the covariance update.
  \item Know why the Joseph form is used and how to implement the update without explicit inverses.
  \item Discretize a continuous-time linear stochastic system exactly and approximately, including Van
        Loan's method.
\end{itemize}
\end{objectives}

Rudolf Kálmán's 1960 paper~\cite{kalman1960} gave a recursive solution to linear estimation that
needs only the current estimate and covariance. It flew on Apollo and remains the backbone of
navigation. This chapter derives it for linear systems. \Cref{ch:eskf} extends it to the nonlinear,
manifold-valued state of VIO.

% ------------------------------------------------------------------------------
\section{Linear-Gaussian state-space models}

A discrete-time linear system with Gaussian noise is described by
\begin{align}
  x_k &= \Phi_{k}\, x_{k-1} + w_k, & w_k &\sim \N(0, Q_k), \label{eq:lin-process}\\
  z_k &= H_k\, x_k + n_k, & n_k &\sim \N(0, R_k), \label{eq:lin-meas}
\end{align}
with $x_0 \sim \N(\hat x_0, P_0)$, and with $x_0$, $\{w_k\}$ and $\{n_k\}$ mutually independent and white.
$\Phi_k$ is the state transition matrix, $Q_k$ the process noise covariance, $H_k$ the measurement
matrix and $R_k$ the measurement noise covariance. The goal is the posterior $p(x_k \mid z_{1:k})$
given all measurements up to time $k$.

% ------------------------------------------------------------------------------
\section{The Bayes filter}

Because the state is Markov (given $x_{k-1}$, the past does not influence $x_k$) and each measurement
depends only on the current state, the posterior can be computed recursively in two steps:
\begin{align}
  \textbf{predict:}&\quad p(x_k \mid z_{1:k-1}) = \int p(x_k \mid x_{k-1})\, p(x_{k-1}\mid z_{1:k-1})\,\dd x_{k-1}, \\
  \textbf{update:}&\quad p(x_k \mid z_{1:k}) \propto p(z_k \mid x_k)\, p(x_k \mid z_{1:k-1}).
\end{align}
For a linear-Gaussian model all of these densities are Gaussian, so each step reduces to updating a mean
and a covariance. We write $\hat x_k^-$, $P_k^-$ for the predicted (prior) moments and $\hat x_k^+$, $P_k^+$ for
the updated (posterior) moments.

% ------------------------------------------------------------------------------
\section{Derivation}

\subsection{Prediction}

By \cref{prop:affine} applied to \eqref{eq:lin-process}, with $x_{k-1}$ and $w_k$ independent,
\begin{equation}
  \hat x_k^- = \Phi_k\,\hat x_{k-1}^+, \qquad P_k^- = \Phi_k P_{k-1}^+\Phi_k\T + Q_k .
  \label{eq:kf-predict}
\end{equation}
Prediction never decreases uncertainty in the directions excited by $Q_k$: it moves the error with the
dynamics and adds process noise.

\subsection{Update}

Before the measurement arrives, the state and the measurement are jointly Gaussian. From
\eqref{eq:lin-meas}, dropping the time index,
\begin{equation}
  \begin{bmatrix} x \\ z \end{bmatrix} \sim \N\!\left(
    \begin{bmatrix} \hat x^- \\ H\hat x^- \end{bmatrix},
    \begin{bmatrix} P^- & P^-H\T \\ HP^- & HP^-H\T + R \end{bmatrix} \right).
\end{equation}
Conditioning on the observed $z$ with \cref{thm:gauss-cond} gives the update directly.

\begin{theorem}[Kalman update]
\label{thm:kf-update}
Define the \emph{innovation} (residual) $r = z - H\hat x^-$, its covariance $S = HP^-H\T + R$, and the
\emph{Kalman gain} $K = P^-H\T S\inv$. Then
\begin{equation}
  \hat x^+ = \hat x^- + K r, \qquad P^+ = P^- - K S K\T = (I - KH)\,P^- .
  \label{eq:kf-update}
\end{equation}
\end{theorem}
\begin{proof}
\Cref{eq:gauss-cond} with $\Sigma_{ab} = P^-H\T$ and $\Sigma_{bb} = S$ gives
$\hat x^+ = \hat x^- + P^-H\T S\inv(z - H\hat x^-)$ and $P^+ = P^- - P^-H\T S\inv HP^-$. Since
$KSK\T = P^-H\T S\inv HP^- = KHP^-$, both covariance forms agree.
\end{proof}

\begin{keyidea}[the gain is a ratio of uncertainties]
For a scalar state measured directly ($H = 1$), $K = \frac{P^-}{P^- + R}$ and
$P^+ = \frac{P^- R}{P^- + R}$. If the prior is much more uncertain than the sensor ($P^- \gg R$), $K \to 1$
and the estimate jumps to the measurement. If the sensor is much worse ($R \gg P^-$), $K \to 0$ and the
measurement is nearly ignored. The posterior variance is always smaller than both. In general, $K$
maps the innovation, which lives in measurement space, into a correction in state space, including
corrections to states that are not measured directly but are correlated with measured ones through
$P^-$. This is how a pose measurement corrects velocity and biases.
\end{keyidea}

\subsection{The information form}

Applying the Woodbury identity (\cref{thm:woodbury}) to $P^+ = P^- - P^-H\T(HP^-H\T + R)\inv HP^-$ gives
\begin{equation}
  (P^+)\inv = (P^-)\inv + H\T R\inv H, \qquad K = P^+ H\T R\inv .
  \label{eq:kf-info}
\end{equation}
Information adds: each independent measurement contributes $H\T R\inv H$. This is the form in which the
MAP problem \eqref{eq:map-ls} is solved, which confirms that the Kalman update is its closed-form
solution. The covariance form \eqref{eq:kf-update} inverts an $m\times m$ matrix ($m$ = measurement
dimension); the information form inverts an $n\times n$ matrix ($n$ = state dimension). Filters use the
former when $m < n$.

% ------------------------------------------------------------------------------
\section{Properties}

\begin{proposition}[Optimality]
For the linear-Gaussian model, $\hat x_k^+ = \E[x_k\mid z_{1:k}]$: the Kalman filter is the minimum
mean-square error estimator among \emph{all} estimators, linear or not. If the noises are not
Gaussian but have the stated means and covariances, it remains the best \emph{linear} unbiased
estimator.
\end{proposition}

\begin{proposition}[Innovation properties]
\label{prop:innovation}
If the model is correct, the innovations $r_k$ are zero-mean, Gaussian, with covariance $S_k$, and
uncorrelated across time. Consequently $r_k\T S_k\inv r_k \sim \chi^2_{m_k}$.
\end{proposition}
\begin{proof}[Sketch]
$r_k = H_k(x_k - \hat x_k^-) + n_k$ is a linear combination of independent zero-mean Gaussians, with
covariance $H_kP_k^-H_k\T + R_k = S_k$. Whiteness follows from the orthogonality principle: the
estimation error is uncorrelated with all past measurements, hence with all past innovations.
The $\chi^2$ statement is \cref{thm:mahalanobis-chi2}.
\end{proof}

\Cref{prop:innovation} is the most useful diagnostic in filtering, because it can be checked on real data
without ground truth. The \emph{normalized innovation squared} (NIS) $r\T S\inv r$ should look like a
$\chi^2_m$ sample. Values that are systematically too large mean the filter is overconfident or the model
is wrong; systematically too small means it is underconfident. Gating (\cref{ch:updates}) uses the same
statistic one measurement at a time.

% ------------------------------------------------------------------------------
\section{The Joseph form}
\label{sec:joseph}

\begin{proposition}[Joseph form]
\label{prop:joseph}
For \emph{any} gain $K$ (optimal or not), the error covariance after the update
$\hat x^+ = \hat x^- + K(z - H\hat x^-)$ is
\begin{equation}
  P^+ = (I - KH)\,P^-\,(I - KH)\T + K R K\T .
  \label{eq:joseph}
\end{equation}
For the optimal gain it equals $(I - KH)P^-$.
\end{proposition}
\begin{proof}
The posterior error is $e^+ = x - \hat x^+ = (I - KH)(x - \hat x^-) - Kn$. Its two terms are independent,
so by \cref{prop:affine} its covariance is \eqref{eq:joseph}. Expanding with $K = P^-H\T S\inv$:
$(I-KH)P^-(I-KH)\T + KRK\T = (I-KH)P^- - (I-KH)P^-H\T K\T + KRK\T = (I-KH)P^- - P^-H\T K\T + K S K\T$,
and the last two terms cancel because $KSK\T = P^-H\T K\T$.
\end{proof}

Why use the longer formula? In floating point, $(I - KH)P^-$ is the difference of two nearly equal
matrices when the update is informative, and it is not symmetric. Rounding can push small eigenvalues
negative, after which the filter is broken: variances become negative, Cholesky fails, gates misbehave.
The Joseph form is a sum of a congruence of $P^-$ and a congruence of $R$, so by \cref{prop:psd} it is PSD
whenever $P^-$ and $R$ are, regardless of rounding in $K$. It also remains correct when $K$ is
deliberately suboptimal (for example, a gain computed with a slightly stale $P$). AGENTS.md F3 mandates
it.

\begin{implnote}[cost of the Joseph form]
With $n$ states and $m$ measurements, $(I - KH)P(I-KH)\T$ costs $\mathcal{O}(n^3)$ if done naively. For the
MSCKF ($n = 75$) this is fine. If it ever matters, note that the update can be rearranged as
$P^+ = P - KHP - (KHP)\T + K S K\T$, which costs $\mathcal{O}(n^2 m)$ and is algebraically identical to the
Joseph form for any gain. It no longer has the sum-of-congruences structure that guarantees positive
semi-definiteness under rounding, though, so keep the symmetrization and monitor the smallest eigenvalue
in tests. Profile before optimizing.
\end{implnote}

% ------------------------------------------------------------------------------
\section{Implementing the update}

\begin{algorithm}[t]
\caption{Numerically careful Kalman update}
\label{alg:kf-update}
\begin{algorithmic}[1]
\Require prior $(\hat x, P)$, measurement $z$, model $(H, R)$
\State $r \gets z - H\hat x$ \Comment{innovation}
\State $S \gets HPH\T + R$; \quad $S \gets \tfrac12(S + S\T)$
\State $L \gets \mathrm{chol}(S)$ \Comment{or LDL$^\mathsf{T}$; failure means $S$ is not PD}
\State $\gamma \gets \|L\inv r\|^2$ \Comment{NIS, for gating and logging}
\If{$\gamma > \chi^2_{0.95}(\dim r)$} \State \Return reject (and count it) \EndIf
\State solve $S\,K\T = HP$ for $K\T$ using $L$ \Comment{never form $S\inv$}
\State $\hat x \gets \hat x + Kr$
\State $P \gets (I - KH)P(I - KH)\T + KRK\T$; \quad $P \gets \tfrac12(P + P\T)$
\end{algorithmic}
\end{algorithm}

\Cref{alg:kf-update} collects the numerical rules of the previous sections. Two details deserve
comment. The gain is obtained by solving $SK\T = HP\T = HP$ (since $P$ is symmetric), so the Cholesky factor
of $S$ is reused for both the gate and the gain. And symmetrization after each step removes the
asymmetry that rounding introduces, which otherwise slowly accumulates (AGENTS.md F2).

\paragraph{Sequential updates.} If $R$ is block-diagonal, the blocks of a stacked measurement are
independent given $x$, and processing them one after another with intermediate updates gives exactly
the same result as a single batch update. This can reduce cost (several small $S$ inversions instead of
one large one) and allows per-block gating. If the blocks share errors (correlated $R$), sequential
processing without accounting for the correlation is wrong.

% ------------------------------------------------------------------------------
\section{From continuous to discrete time}
\label{sec:discretization}

Physical systems, including the error dynamics of an inertial navigation system, are naturally written in
continuous time:
\begin{equation}
  \dot x(t) = F_c\, x(t) + G_c\, w(t), \qquad \E[w(t)w(s)\T] = Q_c\,\delta(t - s),
  \label{eq:ct-model}
\end{equation}
where $w$ is white noise with power spectral density $Q_c$ (\cref{sec:ct-noise}). Over a step of length
$\Dt$ with $F_c$, $G_c$ constant, the solution is
\begin{equation}
  x(t + \Dt) = \underbrace{\exp(F_c\Dt)}_{\Phi}\, x(t) + \underbrace{\int_0^\Dt \exp(F_c s)\,G_c\,w(t + \Dt - s)\,\dd s}_{w_k},
\end{equation}
which is exactly of the form \eqref{eq:lin-process}. The discrete process noise is zero-mean Gaussian with
covariance (by the same white-noise calculation as in \cref{prop:integrated-noise})
\begin{equation}
  Q_d = \int_0^\Dt \exp(F_c s)\, G_c Q_c G_c\T \exp(F_c s)\T\,\dd s .
  \label{eq:qd-exact}
\end{equation}

\subsection{Approximations}

Expanding to first order in $\Dt$ gives the approximations used by most real-time filters and frozen in
AGENTS.md F2:
\begin{equation}
  \Phi \approx I + F_c\Dt, \qquad Q_d \approx G_c Q_c G_c\T\,\Dt .
  \label{eq:disc-first-order}
\end{equation}
Both are accurate to $O(\Dt^2)$ per step. The units confirm \eqref{eq:disc-first-order}: a noise density
$\sigma_a$ in $\mathrm{m/s^2/\sqrt{Hz}}$ gives $\sigma_a^2\Dt$ in $\mathrm{m^2/s^2}$, a velocity variance. A
second-order option adds $\tfrac12F_c^2\Dt^2$ to $\Phi$ and
$\tfrac12(F_cGQG\T + GQG\T F_c\T)\Dt^2$ to $Q_d$.

\begin{example}[Double integrator]
\label{ex:double-integrator-qd}
For position and velocity driven by white acceleration of density $\sigma$, $F_c = \begin{bmatrix}0&1\\0&0\end{bmatrix}$,
$G_c = \begin{bmatrix}0\\1\end{bmatrix}$, $Q_c = \sigma^2$. \Cref{ex:double-integrator} gives $\exp(F_cs)G_c = [s;\ 1]$, so
\eqref{eq:qd-exact} evaluates exactly to
\begin{equation}
  Q_d = \sigma^2\begin{bmatrix} \Dt^3/3 & \Dt^2/2 \\ \Dt^2/2 & \Dt \end{bmatrix}.
\end{equation}
The first-order approximation $Q_d \approx \diag(0, \sigma^2\Dt)$ puts no noise on position within a single
step. Yet after $n$ steps of exact $\Phi$, the accumulated position variance is
$\sum_{k=1}^{n} (k\Dt)^2\sigma^2\Dt \to \sigma^2 t^3/3$ as $\Dt\to0$ (a Riemann sum), which is the exact
answer of \cref{prop:integrated-noise}. The per-step error of the approximation does not accumulate into
a wrong growth law; it only introduces an $O(\Dt)$ relative error, negligible at 200~Hz.
\end{example}

\subsection{Van Loan's method}

When higher accuracy is wanted, \eqref{eq:qd-exact} can be computed exactly with a single matrix
exponential~\cite{vanloan1978}.

\begin{proposition}[Van Loan]
\label{prop:vanloan}
Let $C = \exp\!\Big(\begin{bmatrix} -F_c & G_cQ_cG_c\T \\ 0 & F_c\T \end{bmatrix}\Dt\Big)
= \begin{bmatrix} C_{11} & C_{12} \\ 0 & C_{22}\end{bmatrix}$. Then $\Phi = C_{22}\T$ and $Q_d = \Phi\, C_{12}$.
\end{proposition}
\begin{proof}
Block-triangular exponentials satisfy $C_{22} = \exp(F_c\T\Dt)$ and
$C_{12} = \int_0^\Dt \exp(-F_c(\Dt - s))\,GQG\T\exp(F_c\T s)\,\dd s$ (verify by differentiating in $\Dt$).
Multiplying by $\Phi = \exp(F_c\Dt)$ on the left and substituting $s\to\Dt-s$ gives \eqref{eq:qd-exact}.
\end{proof}

For the 15-state INS this costs a $30\times30$ matrix exponential per step, too expensive at 200~Hz but
excellent as a test oracle for the closed forms of \cref{ch:errordyn}.

% ------------------------------------------------------------------------------
\section{A worked example}

\begin{figure}[t]
  \centering
  \includegraphics[width=\textwidth]{fig_kf1d.pdf}
  \caption{A constant-velocity Kalman filter tracking a 1D target from noisy position measurements.
  (a) The estimate smooths the measurements. (b) The errors stay inside the $\pm3\sigma$ envelopes, which
  shrink quickly from the uninformative prior and settle at the steady-state values set by the balance
  between process noise and measurement noise. The velocity is never measured, yet it is estimated
  through its correlation with position.}
  \label{fig:kf1d}
\end{figure}

\Cref{fig:kf1d} shows the filter on the double-integrator model of \cref{ex:double-integrator-qd}, with
position measurements at 10~Hz. Three behaviors generalize to VIO. The covariance converges from its
initial value to a steady state determined only by $Q$, $R$ and the geometry, not by the data: the
\emph{discrete algebraic Riccati equation}. The unmeasured state (velocity) is estimated through the
correlation that the dynamics build up between position and velocity. And the errors stay inside the
$3\sigma$ envelope about 99.7\% of the time, the visual signature of a consistent filter.

\begin{takeaways}
\begin{itemize}
  \item Predict: $\hat x \gets \Phi\hat x$, $P \gets \Phi P\Phi\T + Q$. Update: condition the joint Gaussian of
        state and measurement.
  \item $K = PH\T S\inv$ maps innovations into state corrections, including corrections to unmeasured but
        correlated states.
  \item Innovations of a correct filter are white with covariance $S$; their NIS is $\chi^2$.
  \item Use the Joseph form, solve rather than invert, and symmetrize.
  \item Discretize with $\Phi \approx I + F_c\Dt$ and $Q_d \approx G_cQ_cG_c\T\Dt$. Van Loan's method gives the exact
        values for testing.
\end{itemize}
\end{takeaways}

\begin{exercises}
  \item Derive \eqref{eq:kf-update} for a scalar state measured twice by independent sensors of variances
        $R_1$ and $R_2$, first sequentially, then as a stacked 2D measurement. Show the results agree.
  \item Show that $P^+ \preceq P^-$ for the optimal gain. Is this still true for an arbitrary gain in the Joseph
        form?
  \item Construct a $2\times2$ example in which $(I - KH)P^-$ computed in single precision is not positive
        semi-definite while the Joseph form is.
  \item Verify \cref{prop:vanloan} numerically for the double integrator, and compare $Q_d$ with the first-order
        approximation for $\Dt = 0.005$ and $\Dt = 0.5$.
  \item For the scalar random walk $x_k = x_{k-1} + w_k$ with $\Var(w) = q$ and direct measurements with
        variance $r$, solve the steady-state Riccati equation for $P^-_\infty$ and the steady-state gain.
  \item Implement \cref{alg:kf-update} in Python, reproduce \cref{fig:kf1d}, and compute the fraction of
        time-steps where the position error lies outside $\pm3\sigma$.
\end{exercises}
